%% ==============================================================================
%% content/02-etat-de-l-art.tex
%% Chapitre 2 : La Classe Racine et les Profils d'Options (.cls & .clo)
%% ==============================================================================
\chapter{La Classe Racine et les Profils d'Options (.cls \& .clo)}
\label{chap:classe-et-options}

\section{Pourquoi Bâtir sur KOMA-Script \code{scrbook} ?}
\label{sec:pourquoi-scrbook}

La distribution standard de \LaTeX{} propose historiquement la classe \code{book}.
Conçue il y a plus de quarante ans, cette classe repose sur des normes typographiques américaines rigides,
avec des marges étroites et un système d'en-têtes peu personnalisable.
Pour modifier son apparence, les utilisateurs sont généralement contraints d'empiler des extensions disparates :
\code{geometry} pour les marges,
\code{fancyhdr} pour les en-têtes,
et \code{titlesec} pour les titres de section.
Or, ces extensions entrent régulièrement en conflit et surchargent les mécanismes internes du moteur.

À l'inverse, la suite KOMA-Script développée par Markus Kohm \autocite{kohm2024koma}
constitue une refonte moderne et cohérente de la mise en page éditoriale.
Fondée sur les règles de proportions harmonieuses de la typographie classique européenne \autocite{bringhurst2004elements,mittelbach2023latex},
la classe \code{scrbook} intègre en son c\oe{}ur :
\begin{itemize}
  \item un calcul dynamique et optique du bloc de texte (\textit{Satzspiegel}) via le paramètre de division \code{DIV},
  \item la prise en compte mathématique de la compensation de reliure (\code{BCOR}),
  \item un système de couches d'en-têtes puissant et natif via \code{scrlayer-scrpage},
  \item et des commandes directes d'ajustement des polices (\code{\textbackslash setkomafont}).
\end{itemize}

\section{Anatomie de la Classe Personnalisée : \code{mythesis.cls}}
\label{sec:anatomie-cls}

Notre classe racine \code{mythesis.cls} encapsule ces fonctionnalités
pour offrir une interface claire et épurée.
Elle prend en charge :

\begin{enumerate}
  \item \textbf{L'interception des options :}
  l'utilisateur peut spécifier \code{[screen]} ou \code{[print]}.
  Toute autre option non reconnue (par exemple la taille de corps \code{12pt} ou l'option \code{draft})
  est automatiquement transmise à la classe parente \code{scrbook} via la macro standard \code{\textbackslash PassOptionsToClass}.
  \item \textbf{La déclaration des métadonnées du document :}
  la classe fournit des accesseurs sémantiques (\code{\textbackslash institution}, \code{\textbackslash degree}, \code{\textbackslash supervisor}, \code{\textbackslash location}, \code{\textbackslash keywords})
  qui alimentent automatiquement la page de garde et les propriétés PDF du visualiseur.
  \item \textbf{La hiérarchie typographique KOMA-Script :}
  les styles des titres sont définis globalement avec \code{\textbackslash setkomafont\{chapter\}} et \code{\textbackslash setkomafont{disposition}},
  garantissant une cohérence visuelle sans modifier le corps du texte.
  \item \textbf{Le chargement conditionnel de l'option de classe (\code{.clo}) :}
  selon le mode choisi, la classe injecte soit \code{mythesis-screen.clo}, soit \code{mythesis-print.clo}.
\end{enumerate}

\section{Le Rôle Déterminant des Fichiers \code{.clo}}
\label{sec:mecanisme-clo}

En \LaTeX{}, l'extension \code{.clo} signifie \textit{Class Option}.
C'est exactement le mécanisme qu'utilise le système \LaTeX{} pour charger les tailles de police de base
(par exemple \code{size10.clo} ou \code{scrsize11pt.clo}).
Dans notre projet, nous détournons élégamment cette fonctionnalité standard
pour définir les profils d'affichage sans alourdir le fichier principal.

\subsection{Le profil écran : \code{mythesis-screen.clo}}
Ce profil est pensé pour la lecture numérique interactive (sur écran d'ordinateur, tablette ou smartphone) :
\begin{itemize}
  \item \textbf{Marges symétriques :} réglage de \code{twoside=false} et \code{BCOR=0mm}.
  Le texte est parfaitement centré sur la page, ce qui offre un confort maximal lors du défilement vertical.
  \item \textbf{Palette chromatique :} couleurs sémantiques (\code{ThesisPrimary},
  \code{ThesisSecondary} et \code{ThesisAccent})
  avec des nuances bleu marine et sarcelle pour un contraste optimal.
  \item \textbf{Hyperliens interactifs :} configuration de \code{hyperref} avec \code{colorlinks=true},
  permettant de distinguer immédiatement les renvois internes (en bleu pétrole), les citations bibliographiques (en vert forêt)
  et les adresses web externes (en violet prune).
  \item \textbf{Signets PDF ouverts :} les panneaux de navigation du lecteur PDF s'ouvrent automatiquement dès l'ouverture du document.
\end{itemize}

\subsection{Le profil impression : \code{mythesis-print.clo}}
Ce profil est optimisé pour la fabrication physique d'un livre ou d'un mémoire relié :
\begin{itemize}
  \item \textbf{Mise en page recto-verso :} réglage de \code{twoside=true} et \code{open=right}.
  Les chapitres majeurs s'ouvrent obligatoirement sur une « belle page » (page impaire située à droite dans le livre ouvert).
  \item \textbf{Compensation de pliure (\code{BCOR}) :} application d'un décalage physique (\code{BCOR=10mm}).
  L'empagement intérieur est décalé pour tenir compte de l'espace absorbé par l'encollage ou la couture de la reliure,
  évitant que le texte ne disparaisse au fond de la gorge du livre.
  \item \textbf{Neutralisation des liens (\code{hidelinks}) :} sur papier, les liens colorés produiraient des trames grisâtres inélégantes.
  L'option \code{hidelinks=true} conserve les liens cliquables dans le PDF numérique sous-jacent,
  mais affiche les textes en noir franc pour une impression laser ou offset impeccable.
\end{itemize}

\section{Tableau Récapitulatif des Deux Profils}
\label{sec:tableau-profils}

Le \cref{tab:comparaison-profils} résume les divergences de comportement entre les deux profils.

\begin{table}[htbp]
  \centering
  \caption{Spécifications comparatives des profils d'options \code{.clo}}
  \label{tab:comparaison-profils}
  \begin{tabularx}{\textwidth}{@{} l >{\raggedright\arraybackslash}X >{\raggedright\arraybackslash}X @{}}
    \toprule
    \textbf{Caractéristique} & \textbf{Profil Écran (\code{mythesis-screen.clo})} & \textbf{Profil Impression (\code{mythesis-print.clo})} \\
    \midrule
    Mode de page           & Recto seul (\code{twoside=false}) & Recto-verso véritable (\code{twoside=true}) \\
    Départ des chapitres   & Page suivante quelconque & Page impaire obligatoire (\code{open=right}) \\
    Compensation reliure   & \qty{0}{\milli\meter} (\code{BCOR=0mm}) & \qty{10}{\milli\meter} (\code{BCOR=10mm}) \\
    Diviseur de grille     & \code{DIV=12} & \code{DIV=12} recalculé avec BCOR \\
    Couleur des liens      & Palette interactive contrastée & Monochrome noir (\code{hidelinks=true}) \\
    En-têtes et folios     & Symétriques (\code{\textbackslash ihead}, \code{\textbackslash ohead}) & Alternés versos/rectos (\code{\textbackslash lehead}, \code{\textbackslash rohead}) \\
    Usage privilégié       & Lecture PDF, soumission en ligne & Reliure thermique, spirale ou cousue \\
    \bottomrule
  \end{tabularx}
\end{table}

\begin{callout}[Recette : Créer un nouveau profil de classe]
Pour ajouter une nouvelle variante (par exemple un profil brouillon \code{mythesis-draft.clo} ou un profil grand format) :
\begin{enumerate}
  \item Créez un fichier \code{mythesis-nom.clo} à la racine.
  \item Déclarez l'option dans \code{mythesis.cls} via \code{\textbackslash DeclareOption\{nom\}\{...\}}.
  \item Activez-le simplement dans \code{main.tex} avec \code{\textbackslash documentclass[nom]\{mythesis\}}.
\end{enumerate}
\end{callout}
